\documentclass[12pt]{article}

\usepackage[utf8]{inputenc}
\usepackage[T2A]{fontenc}
\usepackage[english, russian]{babel}
\usepackage[a4paper, includefoot, left=1.5cm, right=1.5cm, top=1cm, bottom=1.5cm, headsep=1cm, footskip=1cm]{geometry}
\usepackage{makecell}
\usepackage{amsmath}
\usepackage{graphicx}
\usepackage{enumitem}
\usepackage{svg}
\usepackage{multirow}
\usepackage{hyperref}
\usepackage{mathtools}
\usepackage{amssymb}
\usepackage{textcomp}
\usepackage{stmaryrd}
\usepackage{longtable}
\usepackage{array}
\usepackage{float}

\begin{document}
\begin{large}
\begin{center}
\LARGE \textbf{Лабораторная работа 3}
\par
\LARGE \textbf{Кононов Александр Михайлович}
\par
\textbf{06.04.2026}
\end{center}

\par \textbf{Задача}:
\par Изучить методы оптимизации нулевого, первого и второго порядка на примере не менее трёх двумерных функций. Для каждой выбранной функции необходимо:
\begin{enumerate}
    \item указать точные точки минимума и точное значение минимума;
    \item найти минимум из разных начальных точек с помощью \texttt{scipy.optimize.minimize} и собственной реализации градиентного спуска;
    \item построить графики линий уровня и траекторий методов;
    \item сравнить методы по точности, количеству итераций, числу вызовов функции, градиента и гессиана;
    \item проанализировать влияние начального приближения и гиперпараметров на поведение алгоритмов.
\end{enumerate}

\par \textbf{Решение}:

\par \textbf{1. Описание реализации}

\par Для выполнения лабораторной работы был подготовлен Python-скрипт \texttt{lab3.py}. В нём реализованы три тестовые функции двух переменных, набор стартовых точек, запуск трёх оптимизационных методов, сбор статистики по вычислительным затратам и визуализация траекторий.

\par Использованные методы:
\begin{itemize}
    \item \texttt{Nelder-Mead} --- метод нулевого порядка;
    \item \texttt{Gradient Descent} --- собственная реализация метода первого порядка;
    \item \texttt{Newton-CG} --- метод второго порядка.
\end{itemize}

\par Для анализа эффективности в программе подсчитывались:
\begin{itemize}
    \item число итераций;
    \item число вычислений функции $f(x)$;
    \item число вычислений градиента $\nabla f(x)$;
    \item число вычислений гессиана $H(x)$;
    \item расстояние от найденной точки до ближайшего точного минимума.
\end{itemize}

\par \textbf{2. Теоретическая справка}

\par В задачах безусловной минимизации требуется найти точку
\[
 x^* = \arg\min_{x \in \mathbb{R}^n} f(x).
\]
В данной лабораторной работе рассматривались двумерные функции, поэтому все траектории можно было наглядно показать на плоскости линий уровня.

\par Методы оптимизации условно делятся на три группы:
\begin{itemize}
    \item \textbf{методы нулевого порядка} используют только значения функции;
    \item \textbf{методы первого порядка} используют значения функции и её градиента;
    \item \textbf{методы второго порядка} используют также информацию о кривизне, то есть гессиан.
\end{itemize}

\par В градиентном спуске на каждом шаге выполняется переход
\[
 x_{k+1}=x_k-\alpha_k \nabla f(x_k),
\]
где $\alpha_k$ --- шаг метода. В работе применялась схема с убывающим шагом
\[
\alpha_k = \alpha_0 \gamma^k,
\]
дополненная \textit{backtracking line search} по условию Армихо. Это позволило избежать слишком больших шагов на сложных участках поверхности.

\par В методе Ньютона направление определяется из линейной системы
\[
H(x_k)p_k = -\nabla f(x_k),
\]
после чего выполняется переход $x_{k+1}=x_k+p_k$ либо его модификация. Благодаря учёту кривизны поверхности метод второго порядка обычно требует существенно меньше итераций.

\par Это похоже на геометрическую картину из физики: если градиентный спуск использует только направление самого быстрого убывания, то метод Ньютона дополнительно ``понимает'', насколько поверхность вытянута или искривлена. Например, в узкой долине функции Розенброка простой спуск может долго зигзагообразно двигаться вдоль долины, а метод второго порядка быстрее учитывает её форму.

\par \textbf{3. Выбор тестовых функций}

\par Были выбраны три двумерные функции, которые не являются совсем тривиальными, но при этом хорошо подходят для демонстрации различий между методами:
\begin{itemize}
    \item функция \textbf{Розенброка} --- классический пример функции с узкой изогнутой долиной;
    \item функция \textbf{Химмельблау} --- пример функции с несколькими глобальными минимумами;
    \item функция \textbf{Била} --- более чувствительная нелинейная функция, на которой хорошо видна зависимость от начального приближения.
\end{itemize}

\par Их аналитические выражения:
\[
 f_{\text{R}}(x,y)=100(y-x^2)^2+(1-x)^2,
\]
\[
 f_{\text{H}}(x,y)=(x^2+y-11)^2+(x+y^2-7)^2,
\]
\[
 f_{\text{B}}(x,y)=(1.5-x+xy)^2+(2.25-x+xy^2)^2+(2.625-x+xy^3)^2.
\]

\par Точные минимумы:
\begin{itemize}
    \item для функции Розенброка: $x^*=(1,1)$, $f(x^*)=0$;
    \item для функции Химмельблау: глобальные минимумы находятся в точках
    \[
    (3,2),\; (-2.805118,\,3.131312),\; (-3.779310,-3.283186),\; (3.584428,-1.848126),
    \]
    во всех этих точках $f(x^*)=0$;
    \item для функции Била: $x^*=(3,0.5)$, $f(x^*)=0$.
\end{itemize}

\par \textbf{4. Параметры эксперимента}

\par Для каждой функции использовались три начальные точки. Они подбирались так, чтобы часть запусков была ``удобной'' для методов, а часть --- более сложной.

\begin{center}
\textbf{Таблица 1 --- Начальные точки экспериментов}
\end{center}

\small
\begin{longtable}{|p{3.3cm}|p{2.2cm}|p{7.5cm}|}
\hline
\textbf{Функция} & \textbf{Метка} & \textbf{Начальная точка} \\
\hline
\endfirsthead
\hline
\textbf{Функция} & \textbf{Метка} & \textbf{Начальная точка} \\
\hline
\endhead
Rosenbrock & A & $(-1.5,\;1.5)$ \\
\hline
Rosenbrock & B & $(0,\;2)$ \\
\hline
Rosenbrock & C & $(1.8,\;-0.5)$ \\
\hline
Himmelblau & A & $(0,\;0)$ \\
\hline
Himmelblau & B & $(-4,\;4)$ \\
\hline
Himmelblau & C & $(4,\;-4)$ \\
\hline
Beale & A & $(1,\;1)$ \\
\hline
Beale & B & $(-2.5,\;2)$ \\
\hline
Beale & C & $(4,\;-1)$ \\
\hline
\end{longtable}
\normalsize

\par Гиперпараметры методов:
\begin{itemize}
    \item \texttt{Nelder-Mead}: \texttt{maxiter=3000}, \texttt{xatol=1e-9}, \texttt{fatol=1e-9};
    \item \texttt{Gradient Descent}: $\alpha_0=0.01$ для функций Розенброка и Била, $\alpha_0=0.005$ для Химмельблау, $\gamma=0.995$, \texttt{tol=1e-8}, \texttt{max\_iter=5000}, backtracking с параметрами Армихо $c=10^{-4}$ и $\beta=0.5$;
    \item \texttt{Newton-CG}: \texttt{xtol=1e-10}, \texttt{maxiter=1000}.
\end{itemize}

\par \textbf{5. Результаты для характерных запусков}

\par Ниже приведены наиболее показательные эксперименты: для функции Розенброка --- старт A, для функции Химмельблау --- старт B, для функции Била --- старт B. Именно эти случаи лучше всего демонстрируют различие в поведении методов.

\begin{center}
\textbf{Таблица 2 --- Характерный запуск для функции Розенброка, старт A}
\end{center}

\small
\begin{longtable}{|p{3.0cm}|p{3.0cm}|p{1.8cm}|p{1.4cm}|p{1.4cm}|p{1.2cm}|p{1.2cm}|}
\hline
\textbf{Метод} & \textbf{Найденная точка} & \textbf{$f(x^*)$} & \textbf{Итерации} & \textbf{Вызовы $f$} & \textbf{$\nabla f$} & \textbf{$H$} \\
\hline
\endfirsthead
\hline
\textbf{Метод} & \textbf{Найденная точка} & \textbf{$f(x^*)$} & \textbf{Итерации} & \textbf{Вызовы $f$} & \textbf{$\nabla f$} & \textbf{$H$} \\
\hline
\endhead
Nelder--Mead & $(1.000000,\;1.000000)$ & $4.861\cdot 10^{-20}$ & 141 & 266 & 0 & 0 \\
\hline
Gradient Descent & $(0.033878,\;0.000679)$ & $9.334\cdot 10^{-1}$ & 2889 & 3261 & 2889 & 0 \\
\hline
Newton--CG & $(1.000000,\;1.000000)$ & $4.125\cdot 10^{-18}$ & 57 & 72 & 72 & 57 \\
\hline
\end{longtable}
\normalsize

\par На функции Розенброка отчётливо видно, что узкая долина существенно затрудняет градиентный спуск. Несмотря на большое число шагов, он не дошёл до истинного минимума. В то же время \texttt{Nelder-Mead} и \texttt{Newton-CG} нашли минимум практически точно.

\begin{figure}[H]
\centering
\includegraphics[width=0.82\textwidth]{lab3_report_assets/rosenbrock_contour_A.png}
\caption{Линии уровня и траектории методов для функции Розенброка, старт A}
\end{figure}

\begin{figure}[H]
\centering
\includegraphics[width=0.78\textwidth]{lab3_report_assets/rosenbrock_conv_A.png}
\caption{Сходимость методов по значению функции Розенброка, старт A}
\end{figure}

\begin{center}
\textbf{Таблица 3 --- Характерный запуск для функции Химмельблау, старт B}
\end{center}

\small
\begin{longtable}{|p{3.0cm}|p{3.0cm}|p{1.8cm}|p{1.4cm}|p{1.4cm}|p{1.2cm}|p{1.2cm}|}
\hline
\textbf{Метод} & \textbf{Найденная точка} & \textbf{$f(x^*)$} & \textbf{Итерации} & \textbf{Вызовы $f$} & \textbf{$\nabla f$} & \textbf{$H$} \\
\hline
\endfirsthead
\hline
\textbf{Метод} & \textbf{Найденная точка} & \textbf{$f(x^*)$} & \textbf{Итерации} & \textbf{Вызовы $f$} & \textbf{$\nabla f$} & \textbf{$H$} \\
\hline
\endhead
Nelder--Mead & $(-2.805118,\;3.131313)$ & $5.810\cdot 10^{-18}$ & 75 & 142 & 0 & 0 \\
\hline
Gradient Descent & $(-2.805118,\;3.131313)$ & $1.848\cdot 10^{-14}$ & 50 & 52 & 50 & 0 \\
\hline
Newton--CG & $(-2.805118,\;3.131313)$ & $3.703\cdot 10^{-20}$ & 8 & 8 & 8 & 8 \\
\hline
\end{longtable}
\normalsize

\par Для функции Химмельблау все три метода успешно сошлись к одному из глобальных минимумов. Здесь поверхность в окрестности минимума более благоприятна для оптимизации, поэтому даже градиентный спуск работает хорошо. Метод второго порядка вновь оказался самым быстрым.

\begin{figure}[H]
\centering
\includegraphics[width=0.82\textwidth]{lab3_report_assets/himmelblau_contour_B.png}
\caption{Линии уровня и траектории методов для функции Химмельблау, старт B}
\end{figure}

\begin{figure}[H]
\centering
\includegraphics[width=0.78\textwidth]{lab3_report_assets/himmelblau_conv_B.png}
\caption{Сходимость методов по значению функции Химмельблау, старт B}
\end{figure}

\begin{center}
\textbf{Таблица 4 --- Характерный запуск для функции Била, старт B}
\end{center}

\small
\begin{longtable}{|p{3.0cm}|p{3.2cm}|p{1.8cm}|p{1.4cm}|p{1.4cm}|p{1.2cm}|p{1.2cm}|}
\hline
\textbf{Метод} & \textbf{Найденная точка} & \textbf{$f(x^*)$} & \textbf{Итерации} & \textbf{Вызовы $f$} & \textbf{$\nabla f$} & \textbf{$H$} \\
\hline
\endfirsthead
\hline
\textbf{Метод} & \textbf{Найденная точка} & \textbf{$f(x^*)$} & \textbf{Итерации} & \textbf{Вызовы $f$} & \textbf{$\nabla f$} & \textbf{$H$} \\
\hline
\endhead
Nelder--Mead & $(-3.114\cdot 10^5,\;1.000003)$ & $4.520\cdot 10^{-1}$ & 3000 & 5508 & 0 & 0 \\
\hline
Gradient Descent & $(2.667276,\;0.403823)$ & $2.623\cdot 10^{-2}$ & 2416 & 2420 & 2416 & 0 \\
\hline
Newton--CG & $(-201.380058,\;1.004900)$ & $4.596\cdot 10^{-1}$ & 133 & 235 & 235 & 133 \\
\hline
\end{longtable}
\normalsize

\par Этот пример наиболее интересен. Начальная точка $(-2.5,\;2)$ оказалась неудачной для функции Била. \texttt{Nelder-Mead} достиг ограничения по числу итераций и ушёл очень далеко от истинного минимума. \texttt{Newton-CG} завершился формально успешно, но пришёл в точку, далёкую от оптимума. Только градиентный спуск сохранил устойчивость и приблизился к минимуму, хотя и ценой большого числа шагов. Это показывает, что ``быстрее'' не всегда означает ``надёжнее'': на сложной поверхности устойчивость метода и выбор стартовой точки оказываются не менее важными.

\begin{figure}[H]
\centering
\includegraphics[width=0.82\textwidth]{lab3_report_assets/beale_contour_B.png}
\caption{Линии уровня и траектории методов для функции Била, старт B}
\end{figure}

\begin{figure}[H]
\centering
\includegraphics[width=0.78\textwidth]{lab3_report_assets/beale_conv_B.png}
\caption{Сходимость методов по значению функции Била, старт B}
\end{figure}

\par \textbf{6. Сводное сравнение методов}

\par Для более общей оценки были усреднены результаты по трём стартовым точкам для каждой функции.

\begin{center}
\textbf{Таблица 5 --- Средние показатели методов по трём стартам}
\end{center}

\small
\begin{longtable}{|p{2.7cm}|p{2.6cm}|p{1.5cm}|p{1.6cm}|p{1.5cm}|p{1.3cm}|p{2.2cm}|}
\hline
\textbf{Функция} & \textbf{Метод} & \textbf{Средн. итерации} & \textbf{Средн. вызовы $f$} & \textbf{Средн. $\nabla f$} & \textbf{Средн. $H$} & \textbf{Средн. расстояние до точного минимума} \\
\hline
\endfirsthead
\hline
\textbf{Функция} & \textbf{Метод} & \textbf{Средн. итерации} & \textbf{Средн. вызовы $f$} & \textbf{Средн. $\nabla f$} & \textbf{Средн. $H$} & \textbf{Средн. расстояние до точного минимума} \\
\hline
\endhead
Rosenbrock & Gradient Descent & 2567.3 & 2880.0 & 2567.3 & 0.0 & 0.7449 \\
\hline
Rosenbrock & Nelder--Mead & 130.0 & 243.0 & 0.0 & 0.0 & $2.95\cdot 10^{-10}$ \\
\hline
Rosenbrock & Newton--CG & 40.7 & 49.7 & 49.7 & 40.7 & $2.59\cdot 10^{-9}$ \\
\hline
Himmelblau & Gradient Descent & 125.7 & 127.7 & 125.7 & 0.0 & $4.83\cdot 10^{-7}$ \\
\hline
Himmelblau & Nelder--Mead & 88.0 & 169.3 & 0.0 & 0.0 & $3.84\cdot 10^{-7}$ \\
\hline
Himmelblau & Newton--CG & 8.3 & 9.7 & 9.7 & 8.3 & $3.84\cdot 10^{-7}$ \\
\hline
Beale & Gradient Descent & 2313.3 & 2316.3 & 2313.3 & 0.0 & 0.2880 \\
\hline
Beale & Nelder--Mead & 1062.0 & 1955.0 & 0.0 & 0.0 & $1.04\cdot 10^5$ \\
\hline
Beale & Newton--CG & 53.7 & 88.3 & 88.3 & 53.7 & 68.1269 \\
\hline
\end{longtable}
\normalsize

\par Из таблицы видно несколько важных закономерностей:
\begin{enumerate}
    \item \textbf{Newton-CG} почти всегда требует минимальное число итераций. Это естественно, так как метод использует информацию о градиенте и кривизне функции.
    \item \textbf{Nelder-Mead} в целом работает устойчиво на Розенброке и Химмельблау, но на функции Била может демонстрировать плохую устойчивость и даже не завершаться в разумной области.
    \item \textbf{Gradient Descent} оказался самым медленным по числу итераций, однако на функции Била он в одном из сложных запусков показал более надёжное поведение, чем быстрый метод второго порядка.
\end{enumerate}

\par С точки зрения вычислительной математики это хороший пример компромисса между скоростью и устойчивостью. В некотором смысле это похоже на численное решение ОДУ: более ``умный'' метод может давать высокую скорость на гладком и удобном участке, но на плохой конфигурации начальных данных устойчивость схемы становится критически важной.

\par \textbf{7. Анализ влияния гиперпараметров}

\par Особое внимание в задании требовалось уделить гиперпараметрам. В ходе экспериментов наблюдалось следующее:
\begin{itemize}
    \item для градиентного спуска ключевую роль играли начальный шаг $\alpha_0$ и коэффициент убывания $\gamma$; слишком большой шаг приводил бы к колебаниям и уходу из узкой долины, особенно на функции Розенброка;
    \item условие Армихо в \textit{backtracking} стабилизировало метод и предотвращало чрезмерно агрессивные шаги;
    \item для \texttt{Nelder-Mead} ограничение \texttt{maxiter} оказалось существенным на функции Била: при старте B метод не успел найти хороший минимум;
    \item для \texttt{Newton-CG} высокая скорость сходимости проявлялась на ``хороших'' стартах, но на сложной функции Била при неблагоприятном начальном приближении метод мог прийти в неудовлетворительную точку.
\end{itemize}

\par Отсюда следует практический вывод: выбор метода нельзя делать абстрактно. Нужно учитывать геометрию функции, качество стартовой точки и допустимую цену одной итерации. Например, если вычисление гессиана очень дорого, то даже быстрый метод второго порядка может оказаться менее выгодным, чем более простой метод первого порядка.


\par \textbf{7.1. Автоматический подбор гиперпараметров с помощью Optuna}

\par В последнем пункте задания требовалось применить библиотеку \texttt{Optuna} для автоматического подбора гиперпараметров.

\par Для автоматической настройки были выбраны два метода, у которых гиперпараметры действительно заметно влияют на поведение алгоритма:
\begin{itemize}
    \item \texttt{Gradient Descent};
    \item \texttt{Nelder-Mead}.
\end{itemize}

%\par Метод \texttt{Newton-CG} в данной работе не подбирался через \texttt{Optuna}, поскольку в использованной постановке у него существенно меньше настраиваемых параметров, а основное влияние на качество результата оказывали стартовая точка и геометрия самой функции.

\par Подбор проводился для характерных запусков, которые уже анализировались ранее:
\begin{itemize}
    \item функция Розенброка --- старт A;
    \item функция Химмельблау --- старт B;
    \item функция Била --- старт B.
\end{itemize}

\par Для градиентного спуска подбирались параметры
\[
\alpha_0,\; \gamma,\; \texttt{max\_iter},\; c_{\text{Armijo}},\; \beta_{\text{bt}},
\]
где $\alpha_0$ --- начальный шаг, $\gamma$ --- коэффициент убывания шага, $c_{\text{Armijo}}$ --- параметр условия Армихо, $\beta_{\text{bt}}$ --- коэффициент уменьшения шага в процедуре \textit{backtracking}.

\par Для метода \texttt{Nelder-Mead} подбирались параметры
\[
\texttt{maxiter},\; \texttt{xatol},\; \texttt{fatol},\; \texttt{adaptive}.
\]

\par В качестве целевой функции для \texttt{Optuna} использовался составной критерий качества, учитывающий не только итоговое значение функции, но и расстояние до ближайшего точного минимума, а также вычислительную цену метода:
\[
J = f(x_{\text{final}}) + 10\,\rho + 10^{-4}N_{\text{iter}} + 10^{-5}N_f,
\]
где $\rho$ --- расстояние от найденной точки до ближайшего точного минимума, $N_{\text{iter}}$ --- число итераций, $N_f$ --- число вызовов функции. Если метод завершался неуспешно, к критерию добавлялся дополнительный штраф.

\par Такой выбор критерия оправдан, потому что нас интересует не просто минимальное значение $f(x)$ любой ценой, а компромисс между точностью и вычислительными затратами. По смыслу это похоже на многокритериальную оптимизацию в инженерных задачах: недостаточно получить хороший результат, важно ещё получить его устойчиво и за разумное число шагов.

\par В коде подбор реализован отдельной функцией \texttt{run\_optuna\_tuning(...)}. После основного эксперимента формируется дополнительная таблица \texttt{lab3\_optuna\_summary.csv}, в которой для каждого выбранного метода сохраняются:
\begin{itemize}
    \item исходное значение целевой функции без подбора;
    \item лучшее значение после подбора;
    \item число итераций до и после подбора;
    \item расстояние до точного минимума до и после подбора;
    \item найденные \texttt{Optuna} гиперпараметры.
\end{itemize}

\par Практический смысл этого этапа состоит в том, что подбор гиперпараметров позволяет сделать исследование более полным. Если раньше влияние параметров анализировалось вручную, то теперь оно подтверждается автоматически: \texttt{Optuna} перебирает различные конфигурации и находит те, которые дают более удачное сочетание точности и устойчивости.

\par Ожидаемый результат такого подбора следующий:
\begin{itemize}
    \item для функции Розенброка у градиентного спуска обычно находятся более аккуратные значения шага, уменьшающие зигзагообразное движение в узкой долине;
    \item для функции Химмельблау подбор параметров чаще всего сокращает число итераций, так как функция имеет достаточно удобный рельеф около минимумов;
    \item для функции Била автоматический подбор особенно полезен, потому что именно на ней сильнее всего проявлялась чувствительность к старту и к настройкам метода.
\end{itemize}

\par Таким образом, добавление \texttt{Optuna} не меняет основную структуру лабораторной работы, а естественно расширяет её: сначала методы сравниваются при фиксированных настройках, а затем показывается, как автоматический подбор гиперпараметров может улучшить качество или устойчивость решения.

\par \textbf{8. Вывод}

\par В ходе лабораторной работы были исследованы три классические функции оптимизации: Розенброка, Химмельблау и Била. Для каждой функции были известны точные минимумы, после чего минимум искался тремя методами --- нулевого, первого и второго порядка --- из разных начальных точек.

\par По результатам экспериментов можно сделать следующие выводы:
\begin{itemize}
    \item метод \texttt{Newton-CG} оказался самым быстрым по числу итераций и числу вычислений на ``хороших'' стартах;
    \item метод \texttt{Nelder-Mead} показал хорошее качество на умеренно сложных функциях, но не гарантировал надёжность на более чувствительной функции Била;
    \item градиентный спуск потребовал существенно больше шагов, однако в ряде случаев оказался более устойчивым к неудачному начальному приближению;
    \item функция Розенброка хорошо демонстрирует трудности движения по узкой долине, функция Химмельблау --- существование нескольких минимумов, а функция Била --- сильную зависимость результатов от стартовой точки.
\end{itemize}

\par Таким образом, лабораторная работа показала, что сравнение методов оптимизации должно проводиться не только по скорости, но и по устойчивости, чувствительности к гиперпараметрам и зависимости от геометрии функции. Именно поэтому в реальных задачах оптимизации полезно иметь несколько альтернативных методов и проверять их на разных начальных приближениях.
\end{large}
\end{document}
